\section{Methodology}
\label{sec:methodology}

\subsection{HAR and HAR-S Models}
The standard HAR model predicts the log realized volatility using daily, weekly, and monthly (22-day) lags:
\begin{equation}
\log(RV_t) = \alpha + \beta_d \log(RV_{t-1}) + \beta_w \log(RV_{t-5:t-1}) + \beta_m \log(RV_{t-22:t-1}) + \varepsilon_t
\end{equation}
The Sentiment-Augmented HAR (HAR-S) linearly extends this by including the lagged exogenous feature vector $X_{t-1}$, which contains FinBERT sentiment, term slope, credit spread, and microstructure variables:
\begin{equation}
\log(RV_t) = \alpha + \beta_d \log(RV_{t-1}) + \beta_w \log(RV_{t-5:t-1}) + \beta_m \log(RV_{t-22:t-1}) + \gamma^\top X_{t-1} + \varepsilon_t
\end{equation}
These models are estimated via Ordinary Least Squares (OLS) utilizing Heteroskedasticity and Autocorrelation Consistent (HAC) standard errors.

\subsection{The Neural-HAR Architecture}
The Neural-HAR model is a hybrid architecture designed to isolate non-linear interactions without sacrificing the robust autoregressive baseline of the HAR model. The model computes a linear OLS HAR prior. In parallel, a residual neural network takes both the HAR lags and the exogenous features $X_{t-1}$ as input to predict the residual error of the linear prior.

The residual branch is an inverse-validation-loss weighted ensemble of two deep learning architectures:
1. \textbf{Gated Residual Network (GRN):} Utilizes Gated Linear Units (GLU) to perform soft feature selection, dampening noisy sentiment variables while amplifying predictive signals.
2. \textbf{FT-Transformer:} Tokenizes numerical features and employs multi-head self-attention to model complex, higher-order cross-feature interactions.

\subsection{Validation and Evaluation Protocol}
To strictly prevent look-ahead bias, we employ an expanding-window walk-forward validation scheme. The baseline HAR and HAR-S models are re-estimated using all available prior data. For the Neural-HAR model, the walk-forward routine utilizes a 21-day step size with an initial minimum training window of 252 days. At each step, Optuna optimizes the hyperparameters (learning rate, weight decay, dropout, network depth) over 40 trials.

Forecasts are evaluated out-of-sample (OOS) using QLIKE, Mean Squared Error (MSE), and Mean Absolute Error (MAE). 
\begin{equation}
\text{QLIKE}(RV_t, \widehat{RV}_t) = \frac{RV_t}{\widehat{RV}_t} - \log\left(\frac{RV_t}{\widehat{RV}_t}\right) - 1
\end{equation}
We apply the Diebold-Mariano (DM) test with the Harvey-Leybourne-Newbold (HLN) small-sample correction, and the Model Confidence Set (MCS) using a block bootstrap ($B=1000$ resamples) at the $\alpha=0.10$ significance level.

\textbf{Important Limitation on Sample Alignment:} Because the Neural-HAR model is computationally intensive and was evaluated over a shorter trailing period, the sample sizes for the out-of-sample evaluation differ. The baseline HAR expanding-window metrics in our preliminary descriptive statistics cover large samples ($N > 2400$). However, the direct DM and MCS comparisons involving the Neural-HAR model (and the corresponding HAR baseline for that specific evaluation) are strictly evaluated on the shortest common trailing window ($N=275$ for BTC, $N=287$ for SPX, $N=276$ for NIFTY), representing approximately one recent trading year that does \textit{not} include the COVID-19 pandemic shock. All relative statistical comparisons in the subsequent sections are computed on this harmonized short window.
